\documentclass[11pt, a4paper]{article}

% Gemeinsame Style-Definitionen (TEXINPUTS muss templates/ enthalten — siehe scripts/build_tex.py)
\input{bfw-style.tex}

\title{\vspace*{-1.5cm}\Huge\bfseries\color{bfwPrimaryDark}Session~8: Präsentieren, Feedback \& Gesamtwiederholung\\[0.4em]
  \large\color{bfwInkSoft}\textnormal{Vortragstechnik, Reflexion und kompakte Wiederholung des gesamten Lernfeldes~1}}
\author{\textsf{IT Lernfeld 1 --- Das Unternehmen und die eigene Rolle im Betrieb beschreiben}}
\date{}

\begin{document}
\maketitle
\thispagestyle{empty}

\vspace{1em}
\begin{merksatz}[title={\bfseries Worum geht es in dieser Session?}]
In Session~8 führen wir die vorbereitete Betriebspräsentation durch, geben uns gegenseitig
konstruktives Feedback und werten die Teamarbeit mit Reflexionsmethoden aus (Kapitel~1.5.3).
Außerdem wiederholen wir strukturiert den gesamten Stoff des Lernfeldes~1 aus den Sessions~1--7,
um uns gezielt auf die Abschlussprüfung vorzubereiten.
\end{merksatz}

\tableofcontents
\vspace{1em}
\hrule
\vspace{1em}

\section{Vortragstechnik: Stimme und Körpersprache}

\subsection{Nonverbale Kommunikation im Vortrag}

Die Wirkung einer Präsentation hängt nicht allein vom Inhalt ab.
Laut der 55-38-7-Regel (nach Mehrabian; Faustregel, nicht aus dem Lehrwerk --- Gültigkeit kontextabhängig)
bestimmen \textbf{55~\% Körpersprache}, \textbf{38~\% Stimme}
und lediglich \textbf{7~\% der gesprochene Inhalt} den ersten Eindruck beim Publikum.
Das bedeutet: Wer sich auf den Inhalt konzentriert und nonverbale Signale vernachlässigt,
verschenkt mehr als die Hälfte seiner Wirkung.

\begin{center}
\begin{tabular}{p{3.5cm} p{9cm}}
\toprule
\textbf{Element} & \textbf{Hinweise für den Vortrag}\\
\midrule
Körperhaltung & Aufrecht, Füße schulterbreit, beide Beine belastet — keine verschränkten Arme,
                keine Hände in den Hosentaschen\\
\addlinespace
Gestik & Unterstützende, ruhige Handbewegungen einsetzen — übertriebene oder nervöse Gestik
         abbauen\\
\addlinespace
Mimik & Freundlich und zugewandt — der Blick in die Folien unterbricht den Kontakt zum Publikum\\
\addlinespace
Blickkontakt & Über den gesamten Raum verteilen, nicht auf eine Person oder die Leinwand fixieren\\
\addlinespace
Lautstärke & Klar und deutlich sprechen — auch die letzte Reihe muss alles verstehen\\
\addlinespace
Tempo & Angemessen langsam — bewusste Pausen einsetzen, um Wichtiges wirken zu lassen\\
\addlinespace
Betonung & Schlüsselbegriffe hervorheben, Monotonie vermeiden\\
\bottomrule
\end{tabular}
\end{center}

\begin{merksatz}[title={\bfseries Merke: 55-38-7-Regel}]
55~\% Körpersprache, 38~\% Stimme, 7~\% Inhalt --- so gewichtet die Kommunikationsforschung
die Wirkfaktoren beim ersten Eindruck. Nonverbale Kompetenz ist daher kein Zusatz,
sondern ein zentraler Bestandteil professioneller Präsentation.
(nach Mehrabian; Faustregel, nicht aus dem Lehrwerk --- Gültigkeit kontextabhängig)
\end{merksatz}

\subsection{Umgang mit Lampenfieber}

Lampenfieber ist eine normale körperliche Reaktion auf eine Stresssituation und kein Zeichen
von Schwäche. Das ausgeschüttete Adrenalin schärft die Konzentration und steigert die
Reaktionsfähigkeit --- leichtes Aufgeregt-Sein ist ein echtes Leistungsplus.

Strategien zur Bewältigung:

\begin{itemize}
  \item \textbf{Gründliche Vorbereitung:} Die wirksamste Maßnahme --- wer den Stoff sicher
        beherrscht und den roten Faden kennt, tritt deutlich ruhiger auf.
  \item \textbf{Probevortrag:} Einen kompletten Vortrag einmal üben --- allein oder vor
        vertrauten Personen. Dabei Zeitnahme und ggf. Videoaufnahme nutzen.
  \item \textbf{Atemübungen:} Kurz vor dem Auftritt tief einatmen (4 Sekunden), Luft kurz
        anhalten, langsam ausatmen (6 Sekunden) --- beruhigt das Nervensystem spürbar.
  \item \textbf{Positive Selbstinstruktion:} An frühere Erfolge denken statt an mögliche Fehler.
        Das Publikum mental als Partner wahrnehmen, der den Vortrag gelingen sehen möchte.
  \item \textbf{Umdeuten:} Lampenfieber als Signal von Verantwortungsgefühl und Einsatz
        interpretieren --- nicht als Schwäche.
\end{itemize}

\section{Feedbackregeln}

\subsection{Was ist konstruktives Feedback?}

Konstruktives Feedback ist ein gezieltes Entwicklungsangebot, das die Leistung der empfangenden
Person stärken soll --- kein Angriff und keine persönliche Kritik.
Gutes Feedback folgt sechs Grundregeln:

\begin{enumerate}
  \item \textbf{Konkret:} Genaue Beobachtungen benennen, keine vagen Eindrücke.
        Nicht: \emph{„Sie wirkten nervös"} --- sondern: \emph{„Sie haben dreimal auf die Folie
        geschaut."}
  \item \textbf{Beschreibend statt wertend:} Was wurde wahrgenommen?
        Keine Urteile über die Person --- nur Beschreibungen des Verhaltens.
  \item \textbf{Als Ich-Botschaft:} \emph{„Ich habe bemerkt, dass \ldots"} --- nicht
        \emph{„Sie haben falsch gemacht \ldots"} Ich-Botschaften vermeiden Schuldzuweisungen.
  \item \textbf{Zeitnah:} Feedback direkt nach der Beobachtung geben ---
        nicht Tage später, wenn die Situation nicht mehr präsent ist.
  \item \textbf{Mit Positivem beginnen:} Stärken zuerst nennen, dann Verbesserungshinweise ---
        das erhält die Motivation und öffnet die Bereitschaft zum Zuhören.
  \item \textbf{Nur auf Veränderbares beziehen:} Hinweise sollen der Person nützen ---
        daher nur auf Verhaltensweisen eingehen, die sie tatsächlich ändern kann.
\end{enumerate}

\subsection{Feedback annehmen}

Auch das Empfangen von Feedback ist eine erlernbare Kompetenz:
Zuhören ohne sofortige Rechtfertigung, Nachfragen bei Unklarheiten,
Feedback als wertvolles Angebot betrachten und selbst entscheiden,
welche Hinweise man annehmen und umsetzen möchte.

\section{Präsentationsbewertung und Reflexion}

\subsection{Bewertungsbogen für Präsentationen}

Ein Bewertungsbogen macht die Beurteilung transparent und vergleichbar.
Er wird \emph{vor} der Präsentation ausgegeben, damit alle Beobachtenden nach denselben
Maßstäben urteilen.

\begin{center}
\begin{tabular}{p{4.5cm} p{1.5cm} p{7.5cm}}
\toprule
\textbf{Kriterium} & \textbf{Gewicht} & \textbf{Beschreibung}\\
\midrule
Fachlicher Inhalt & 40~\% & Richtigkeit, Vollständigkeit, roter Faden, Quellenangaben\\
\addlinespace
Struktur & 20~\% & Einstieg -- Hauptteil -- Schluss erkennbar; Übergänge klar formuliert\\
\addlinespace
Medien und Folien & 20~\% & Übersichtlichkeit, Lesbarkeit ($\geq$30~Punkt), keine Textwüsten,
                              sinnvolle Visualisierungen\\
\addlinespace
Vortragsstil & 20~\% & Körpersprache, Blickkontakt, Stimme, Lampenfieber-Management\\
\midrule
\textbf{Gesamt} & 100~\% & \\
\bottomrule
\end{tabular}
\end{center}

Der Vergleich von \textbf{Selbst- und Fremdbewertung} nach der Präsentation fördert die
Reflexionskompetenz: Blinde Flecken werden sichtbar --- Stärken, die man unterschätzt,
und Schwächen, derer man sich noch nicht bewusst war.

\subsection{Strategien und Methoden zur Reflexion (Kapitel 1.5.3)}

Das Lehrbuch unterscheidet zwei komplementäre Strategien:

\begin{center}
\begin{tabular}{p{3.5cm} p{4.5cm} p{5cm}}
\toprule
\textbf{Strategie} & \textbf{Ziel} & \textbf{Methoden (Beispiele)}\\
\midrule
Reflexions"-strategie & Eigene Arbeit selbstkritisch hinterfragen, Defizite erkennen &
  Blitzlicht, Spinnennetz, Stimmungsbarometer, Lerntagebuch, Fragebogen, Interview\\
\addlinespace
Regulations"-strategie & Konsequenzen ziehen, Defizite beheben, Kompetenzen ausbauen &
  Zielvereinbarungen, Lernlisten, Korrekturen erstellen, Lernstrategie anpassen\\
\bottomrule
\end{tabular}
\end{center}

\textbf{Blitzlicht:} Jede Person nennt in einem Satz, was gut lief und was verbessert werden
sollte. Schnelle Rundrunde --- gut geeignet für kurze Zeitfenster.

\textbf{Spinnennetz:} Die Gruppe bewertet mehrere Kriterien (z.\,B. Zusammenarbeit,
Zeitmanagement, Qualität) auf einer Skala. Die Werte werden in einem Kreisdiagramm eingetragen
und machen Stärken und Entwicklungsfelder sichtbar.

\section{Gesamtwiederholung: Lernfeld~1 im Überblick}

\subsection{Sessions 1 und 2 --- IT-Berufe und Ausbildungsrecht}

\begin{center}
\begin{tabular}{p{3cm} p{10.5cm}}
\toprule
\textbf{Thema} & \textbf{Kernaussagen}\\
\midrule
IT-Berufe & 4 Berufe seit 2020: Fachinformatiker/-in (4 Fachrichtungen), IT-Systemelektroniker/-in,
            Kfm. IT-System-Management, Kfm. Digitalisierungsmanagement\\
\addlinespace
Duales System & Zwei Lernorte: Betrieb (Ausbildungsordnung) und Berufsschule (Rahmenlehrplan);
                Rechtsgrundlage: BBiG\\
\addlinespace
BBiG & § 11: Pflicht zur Vertragsniederschrift; § 20: Probezeit 1--4 Monate;
       § 22: Kündigung (in Probezeit ohne Frist)\\
\addlinespace
JArbSchG & § 8: max. 8~h täglich / 40~h wöchentlich; § 19: Urlaubsstaffelung
           (unter 16~J. 30 Tage, unter 17~J. 27 Tage, unter 18~J. 25 Tage)\\
\addlinespace
JAV & Interessenvertretung der Jugendlichen und Auszubildenden nach § 70~BetrVG\\
\bottomrule
\end{tabular}
\end{center}

\subsection{Sessions 3 und 4 --- Betrieb, Unternehmen, Rechtsformen und Organisation}

\begin{center}
\begin{tabular}{p{3cm} p{10.5cm}}
\toprule
\textbf{Thema} & \textbf{Kernaussagen}\\
\midrule
Begriffe & Betrieb = örtlich-technische Einheit; Unternehmen = wirtschaftliche Einheit;
           Firma = Handelsname nach § 17~HGB\\
\addlinespace
JIKU & Hamburger Standort: 52 Mitarbeiter, 8 Auszubildende; Verbund: 10 Systemhäuser, 480~MA\\
\addlinespace
Ziele & Sach-, Formal- und Sozialziele; SMART-Formel; komplementäre und konkurrierende Ziele\\
\addlinespace
Rechtsformen & UG ab 1~€; GmbH ab 25.000~€; AG ab 50.000~€;
               OHG/KG/GbR: persönliche Haftung, keine Mindestkapitalvorgabe\\
\addlinespace
Organisation & Einlinien-, Mehrliniensystem, Stablinienorganisation, Matrixorganisation;
               Aufbauorganisation vs. Ablauforganisation\\
\bottomrule
\end{tabular}
\end{center}

\subsection{Sessions 5 und 6 --- Wertschöpfung, Kreislauf und Märkte}

\begin{center}
\begin{tabular}{p{3cm} p{10.5cm}}
\toprule
\textbf{Thema} & \textbf{Kernaussagen}\\
\midrule
Wertschöpfung & Formel: Umsatzerlöse $-$ Vorleistungen $-$ Abschreibungen;
                verteilt an Mitarbeiter, Fremdkapitalgeber, Staat, Eigenkapitalgeber\\
\addlinespace
Prozesse & Kernprozesse (direkte Wertschöpfung) vs. Supportprozesse (keine direkte Wertschöpfung);
           IPERKA als Methode der vollständigen Handlung\\
\addlinespace
Produktions"-faktoren & Gutenberg: Elementarfaktoren (Arbeit, Betriebsmittel, Werkstoffe)
                         + dispositiver Faktor (Leitung);
                         VWL: Arbeit, Boden, Kapital (+ Wissen als 4. Faktor)\\
\addlinespace
Güterarten & Freie vs. wirtschaftliche Güter; Konsum- vs. Investitionsgüter;
             Sach- vs. Dienstleistungen\\
\addlinespace
Wirtschafts"-kreislauf & Haushalte $\leftrightarrow$ Unternehmen $\leftrightarrow$ Staat
                          $\leftrightarrow$ Banken $\leftrightarrow$ Ausland\\
\addlinespace
Marktformen & Monopol (1 Anbieter), Oligopol (wenige), Polypol (viele);
              Gleichgewichtspreis am Schnittpunkt von Angebot und Nachfrage\\
\addlinespace
IT-Markt & Netzwerkeffekte, Lock-in-Effekte, Plattformökonomie, Winner-takes-all\\
\bottomrule
\end{tabular}
\end{center}

\subsection{Session 7 --- Präsentation vorbereiten und planen}

\begin{center}
\begin{tabular}{p{3cm} p{10.5cm}}
\toprule
\textbf{Thema} & \textbf{Kernaussagen}\\
\midrule
IPERKA & Informieren -- Planen -- Entscheiden -- Realisieren -- Kontrollieren -- Auswerten\\
\addlinespace
Tuckman & Forming -- Storming -- Norming -- Performing (-- Adjourning)\\
\addlinespace
Planung & Rahmenbedingungen klären; Zielgruppe analysieren; Inhalte nach KISS und
          10-20-30-Regel auswählen; Aufbau: Einstieg -- Hauptteil -- Schluss\\
\addlinespace
IHK-Bewertung & Aufbau/Struktur (10~P.), Sprachliche Gestaltung (10~P.),
                Zielgruppengerechte Darstellung (10~P.) --- gesamt max. 30~Punkte\\
\bottomrule
\end{tabular}
\end{center}

\begin{merksatz}[title={\bfseries Prüfungstipp: Querverbindungen}]
In der Abschlussprüfung werden Themen aus verschiedenen Kapiteln verknüpft.
Beispiel: Rechtsform (GmbH) $\Rightarrow$ beschränkte Haftung $\Rightarrow$ niedriges Risiko
für Gesellschafter $\Rightarrow$ beeinflusst Formalziel und Risikostrategie.
Wer Verbindungen erkennt und benennen kann, zeigt echtes Systemverständnis.
\end{merksatz}

\end{document}
